\documentclass{article}
\usepackage[utf8]{inputenc}
\usepackage{amsmath}
\usepackage{geometry}
\geometry{margin=2cm}
\begin{document}

\section*{Emissões Veiculares e Razão Ar/Combustível}

\subsection*{Interpretação do Enunciado e Estratégia de Resolução}
A questão apresenta um gráfico que relaciona as emissões de diversos gases poluentes (CO_2, HC, CO, NO e O_2) com a razão ar/combustível em motores a gasolina. A tarefa consiste em identificar, para a razão ar/combustível igual a 18, o comportamento das emissões, especificamente de CO e CO_2.

Para isso, devemos:
\begin{itemize}
  \item Localizar a razão 18 no eixo horizontal do gráfico.
  \item Analisar as curvas correspondentes a CO e CO_2 nesse ponto.
  \item Comparar os níveis relativos dessas emissões para concluir se são altas ou baixas.
  \item Confirmar qual alternativa descreve essa situação corretamente.
\end{itemize}

\subsection*{Conteúdo e Mini Revisão Teórica}
\begin{itemize}
  \item \textbf{Razão ar/combustível}: proporção entre a massa de ar e de combustível na mistura. A razão estequiométrica (cerca de 15 para gasolina) corresponde à combustão ideal, com todo combustível totalmente queimado.

  \item \textbf{Combustão completa}: ocorre quando há ar suficiente para oxidar totalmente o combustível. Isso resulta em baixa emissão de gases tóxicos incompletos como CO e HC e alta emissão de CO_2.

  \item \textbf{Emissões veiculares}: gases liberados pela queima do combustível, influenciados pela mistura. Em excesso de ar, CO e HC diminuem e CO_2 e O_2 aumentam.
\end{itemize}

\subsection*{Resolução e Pulo do Gato}

Examinando o gráfico, na razão ar/combustível igual a 18:

\begin{itemize}
  \item \textbf{CO}: Está próximo de zero, indicando baixa emissão.
  \item \textbf{CO_2}: Próximo do pico máximo, indicando alta emissão.
\end{itemize}

Assim, a alternativa que indica baixa emissão de CO e alta de CO_2 (alternativa C) está correta.

\textit{Pulo do gato:} Em mistura com excesso de ar, a combustão é mais completa, reduzindo a emissão de CO (monóxido de carbono) e aumentando a de CO_2 (dióxido de carbono).

\subsection*{Armadilhas e Distratores}

Muitos confundem excesso de ar com combustão incompleta, esperando aumento de CO ou HC. Outros interpretam incorretamente o comportamento do NO e O_2 ou se confundem com os eixos do gráfico remendando respostas erradas.

\subsection*{Código da Questão}
\texttt{CN\_QUIMICA\_ITEM\_001}

\bigskip

\textbf{Fonte do gráfico:} RANGEL, M. C.; CARVALHO, M. F. A. Impacto dos catalisadores automotivos no controle da qualidade do ar. \textit{Química Nova}, v. 26, 2003.

\begin{center}
\textit{\small
Gráfico contendo % volume ou ppm dos gases O_2, CO_2, HC, CO e NO em função da razão ar/combustível (massa/massa), onde a razão estequiométrica 15 é indicada e a região à direita (incluindo 18) representa excesso de ar.}
\end{center}



\end{document}